% 第1章 引言

\section{研究背景}

让机器理解三维室内场景——识别房间里的桌子、椅子、沙发，搞清楚谁在谁旁边、谁在谁上面——是计算机视觉和机器人领域的一个基础问题。这类任务需要标注数据：每张图中每个物体属于什么类别，它们之间的空间关系是怎样的。

人工标注三维场景的成本很高。一个包含约 70 个物体的中等规模房间，从头到尾标完需要几个小时的专家劳动，而且标注者的空间判断还会引入主观误差。视觉语言模型（VLM）可以自动生成描述，算是个省力的替代方案，但生成质量经常不过关。这类模型倾向于输出诸如"自然光洒满整个空间"这样的全局氛围描述，而不是"靠墙放着一张棕色木桌"这种具体的物体描述。Song 等人\cite{llm_planner} 的测试提供了一个量化参考：在 Stanford VLN-CE 长程导航任务中，当步骤超过 5 步时 VLM 的决策准确率下降 37.2\%，环境扰动让动作成功率再降 42.8\%。

Tong 等人\cite{code2logic} 从游戏领域找到了一条有意思的解决路径。他们提出的 Code2Logic 方法基于一个简单的观察：游戏源代码里的坐标数组、碰撞检测逻辑、胜负判定规则本身就是天然的结构化标注。让大语言模型读游戏代码，按模板生成问答对，就能零边际成本地无限产出完美的标注数据。他们构建的 GameQA 数据集包含了 30 款游戏和 14 万条问题，训练后模型在 7 个视觉基准上平均提升了 2.33 个百分点。

这就引出了一个自然的迁移问题：能不能把游戏代码驱动的这套自动化方法搬进真实三维场景？

这个迁移至少面临三个困难。游戏逻辑是显式的，物体在哪、能做什么都由代码精确指定；真实场景的逻辑是隐式的，物体要从 RGB-D 数据里分割和识别。游戏画面是完美渲染的，没有噪声；真实的传感器数据有噪声、遮挡和光照变化。游戏里的每一步操作都可以和代码对比来验对错，真实场景没有这种天然的正确答案。本文的工作围绕这三个困难展开：把 Code2Logic 的模板化方法从 2D 游戏迁移到 3D 场景，构建一条端到端的自动标注管线，再用合成真值系统来验证标注精度。

\section{本文工作}

本文搭建了一条三层数据合成链路。底层是 Code2Logic，30 款游戏的 1397 条 QA 证明了从代码逻辑到模板再到数据这条路可行。中层是 Concept-Graphs，整合了 GradSLAM、Grounded-SAM、CF-SLAM、LLaVA 和 GPU 推理，串成了一条从 RGB-D 图像序列到结构化 QA 数据集的六步流水线。上层是 Isaac Sim，生成 20 个合成场景，每个物体的坐标、类别和尺寸都是精确已知的，用作验证标注精度的金标准。三层递进：方法在底层验证，在中层迁移到真实场景，在上层被合成真值校准。

本文修了一条端到端的自动标注管线。中间六个步骤依次是：GradSLAM 做 3D 重建，Grounded-SAM 做物体分割，CF-SLAM 做三维建图，LLaVA v1.5 为每个物体裁剪区域生成自然语言描述，llama.cpp 在 GPU 上批量优化描述并推理物体间的空间关系，最后按模板生成 QA 数据集。管线的关键改进在第五步的 prompt 设计。LLaVA 的原始描述有两个常见问题：太笼统（"自然光洒满空间"）和输出被截断（"es are all white..."）。重写后的 prompt 要求模型显式输出物体类别（door、table、sofa 等），同时明确禁止了 light、atmosphere、depth 等十几个场景级描述词。这个改动把物体类别识别率从接近零提到了约 85\%。

本文还对 GPU 推理做了一次深度优化。Ollama 自带的 CUDA 12.8 库与服务器驱动 565.57 不兼容，导致推理全部落在 CPU 上，实测吞吐量只有约 3 tok/s。换用源码编译的 llama-cpp-python（链接系统 CUDA 12.6 库），qwen2.5-3B 模型的全部 29 层加载到单张 RTX 4090 显存上。优化后吞吐量达到 99 tok/s，单个场景的 LLM 处理时间从 15 分钟降到了 50 秒。在 5 款国产开源模型上的横向对比表明，Qwen2.5-3B 在速度（86 tok/s）、质量（100\%）和格式合法性（100\%）三个维度上综合最优，显存仅需 1.9 GB。

本文建了一套跨源评估体系。从数据量、推理类型分布、关系类型分布、标注精度和管线效率五个维度，对比了 Code2Logic（游戏合成，1397 QA）、Concept-Graphs（真实 3D 推断，370 QA）和 Isaac Sim（合成 3D 真值，795 QA）三源数据。三类数据累积 2395 条 QA，覆盖五种推理类型。彼此之间是互补关系，不是替代关系：Code2Logic 提供涉及多步规划的 Strategy Optimization 类 QA（这是静态三维场景里不出现的推理类型），Isaac Sim 贡献了大量高精度空间关系（610 条，准确率 100\%），Concept-Graphs 则用开放词汇模型弥补了 Isaac Sim 只能识别预设家具类别的局限。

\section{关键术语与问题边界}

本文中的"标注数据"特指结构化的视觉问答格式，每条包含一个关于场景中物体的问题、一个答案和一个类型标签。选择 VQA 格式是因为它目前是视觉语言模型训练和评估的主流范式，产出的数据可以直接用于下游微调，而且问题-答案的结构天然适合模板化生成。管线的自动化程度是从 RGB-D 数据到 QA 数据集的自动转化，中间不需要人工标注或人工设计规则，但管线依赖公开的 RGB-D 数据集作为输入，大语言模型的 prompt 模板也需要人工设计。标注质量的评测目前依赖 Isaac Sim 合成场景的 ground truth 作为推断标注的验证基准，加上作者对部分 QA 的人工定性抽查。大规模人工评估和下游 VLM 训练验证尚未完成，这两个局限在第五章中有具体讨论。本文方法的设计和验证均限于室内静态场景，室外场景、工业环境和包含动态物体的场景不在当前实验范围内。

\section{论文结构}

第二章梳理相关技术：VLM、CoT 推理、Code2Logic、三维场景图构建和合成数据方法。第三章展开方法细节，包括三层架构、六步管线、prompt 设计、GPU 加速方案和 Isaac Sim 验证系统。第四章报告实验：7 个场景的标注产出、描述质量前后对比、prompt 消融实验、5 模型选型基准、管线效率分析和三源数据多维对比。第五章总结全文，讨论局限性和后续方向。
